\documentclass[11pt]{article}
\usepackage[a4paper,margin=1in]{geometry}
\usepackage{amsmath,amssymb,amsthm,mathtools}
\usepackage{enumitem}
\usepackage[hidelinks]{hyperref}
\usepackage{microtype}

\theoremstyle{definition}
\newtheorem*{theorem*}{Corrected theorem}

\newcommand{\Nzero}{\mathbb N_0}
\newcommand{\dd}{\partial_t}
\newcommand{\Ltwo}{L^2(J)}
\newcommand{\K}{K}
\newcommand{\gam}{\gamma_B}
\newcommand{\Yone}{Y_1}
\newcommand{\Y}{Y}
\newcommand{\V}{V}
\newcommand{\Vprime}{V'}
\newcommand{\hatF}{\widehat F}
\newcommand{\dotE}{\dot E}
\newcommand{\hatY}{\widehat Y}
\newcommand{\clLtwo}[1]{\overline{#1}^{\,L^2(J)}}
\newcommand{\ip}[2]{\langle #1,#2\rangle}

\title{Theorem 5.2.18 errata}
\author{RA}
\date{\today}

\begin{document}

\maketitle

\begin{abstract}
	\noindent
This note records errata to Theorem 5.2.18 in 
\begin{quote}
RA,
\emph{Stability of space-time Petrov-Galerkin discretizations for
parabolic evolution equations}, 2012,
\href{https://doi.org/10.3929/ethz-a-007563932}{doi:10.3929/ethz-a-007563932}.
\end{quote}
The theorem's conclusion remains the uniform discrete inf-sup estimate
\[
  \inf_{L\in\Nzero}\gam(X_L,Y_L)>0.
\]
The corrections below make explicit the intended meaning of the sparse
tensor notation, add the closedness convention needed in the abstract
statement, and fill in the omitted transfer from the auxiliary test space
to the actual test space.
\end{abstract}

\section*{1. Corrections to the surrounding notation}

\begin{enumerate}[label=\textbf{C\arabic*.},leftmargin=2.5em]
\item \textbf{Sparse tensor symbols.}
The symbols \(\bigcup\) in Proposition 4.4.12, Proposition 5.2.14, Remark 5.2.15, and Theorem 5.2.18 should be read as linear sums, not literal set-theoretic unions.  Thus, for example,
\[
  X_L=\sum_{0\le k+\ell\le L}E_k\otimes U_\ell,
  \qquad
  Y_{L,1}=\sum_{0\le k+\ell\le L}F_k\otimes V_\ell .
\]
A literal union of subspaces need not be a subspace.

\item \textbf{Auxiliary temporal spaces.}
In the proof of Theorem 5.2.18 the auxiliary temporal space \(E_k+\dd E_k\) should be replaced, in the infinite-dimensional abstract formulation, by its \(L^2(J)\)-closure:
\[
  \hatF_k:=\clLtwo{E_k+\dd E_k},
  \qquad
  \dotE_k:=\clLtwo{\dd E_k}.
\]
For the finite-dimensional spline and wavelet spaces used later in the thesis, the closures are unnecessary.

\item \textbf{Auxiliary test space.}
The auxiliary test space should be written as
\[
  \hatY_L=\hatY_{L,1}\times V_L,
  \qquad
  \hatY_{L,1}:=\sum_{0\le k+\ell\le L}\hatF_k\otimes V_\ell .
\]
If the printed expression is interpreted as a sum of product subspaces, its second component is \(V_L\) by nestedness of the \(V_\ell\).
\end{enumerate}

\section*{2. Corrected theorem, basic sparse case}

\begin{theorem*}
Let \((E_k,F_k)_{k\in\Nzero}\) be an admissible temporal pair in the sense of Definition 5.2.17.  Let
\[
  U_\ell\subseteq V_\ell\subseteq V,
  \qquad \ell\in\Nzero,
\]
be nested non-trivial closed spatial subspaces satisfying
\[
  \eta:=\inf_{\ell\in\Nzero}\K_{V'\times V}(U_\ell,V_\ell)>0.
\]
Assume, in addition, either that all building blocks are finite-dimensional, or that the auxiliary spaces defined in \textnormal{C2--C3} and the finite sparse tensor sums below are closed subspaces of the ambient Hilbert spaces.

For \(L\in\Nzero\), define
\[
  X_L:=\sum_{0\le k+\ell\le L}E_k\otimes U_\ell,
  \qquad
  Y_{L,1}:=\sum_{0\le k+\ell\le L}F_k\otimes V_\ell,
  \qquad
  Y_L:=Y_{L,1}\times V_L .
\]
Then
\[
  \inf_{L\in\Nzero}\gam(X_L,Y_L)>0.
\]
The full tensor, anisotropic sparse tensor, and finite-sum variants in Remark 5.2.15 remain valid after the same changes.
\end{theorem*}

\section*{3. Proof of the corrected theorem}

\paragraph{Step 1: admissibility passes to the closed auxiliary space.}
Definition 5.2.17 gives
\[
  c_T^0:=\inf_{k\in\Nzero}
  \K_{\Ltwo\times\Ltwo}(E_k+\dd E_k,F_k)>0 .
\]
Let \(P_k:\Ltwo\to F_k\) be the \(L^2(J)\)-orthogonal projector.  Since \(F_k\) is closed,
\[
  \|P_k f\|_{L^2(J)}\ge c_T^0\|f\|_{L^2(J)}
  \qquad\forall f\in E_k+\dd E_k.
\]
By continuity, the same estimate holds for all \(f\in\hatF_k\).  Hence
\[
  c_T:=\inf_{k\in\Nzero}\K_{\Ltwo\times\Ltwo}(\hatF_k,F_k)\ge c_T^0>0 .
\]

\paragraph{Step 2: stability with the auxiliary test space.}
The corrected auxiliary space satisfies the inclusion hypothesis of Corollary 5.2.13:
\[
  X_L\times\{w_L(0):w_L\in X_L\}\subseteq \hatY_L,
\]
because \(E_k\subseteq \hatF_k\) and \(U_\ell\subseteq V_\ell\subseteq V_L\) for \(\ell\le L\).

If \(\dd X_L\ne\{0\}\), apply Proposition 4.4.12 (equivalently, the tensor-product step behind Proposition 5.2.14) to the temporal pair \((\dotE_k,\hatF_k)\), omitting any initial zero derivative levels.  For every nonzero \(\dotE_k\),
\[
  \K_{\Ltwo\times\Ltwo}(\dotE_k,\hatF_k)=1
\]
because \(\dotE_k\subseteq\hatF_k\).  Therefore the tensor-product estimate gives
\[
  \K_{Y_1'\times Y_1}
  \left(
    \sum_{0\le k+\ell\le L}\dotE_k\otimes U_\ell,
    \hatY_{L,1}
  \right)
  \ge \eta .
\]
Since
\[
  \dd X_L\subseteq \sum_{0\le k+\ell\le L}\dotE_k\otimes U_\ell,
\]
we also have
\[
  \K_{Y_1'\times Y_1}(\dd X_L,\hatY_{L,1})\ge \eta .
\]
Corollary 5.2.13 then yields a constant \(c_1>0\), independent of \(L\), such that
\[
  \gam(X_L,\hatY_L)\ge c_1
\]
for all levels with \(\dd X_L\ne\{0\}\).

If \(\dd X_L=\{0\}\), then each \(w\in X_L\) is time-independent.  Testing with \(v=(w,w(0))\in\hatY_L\) and using \(a_{\rm shift}=0\) gives
\[
  \ip{Bw}{v}
  =\int_J a(t;w,w)\,dt+\|w(0)\|_H^2
  \ge a_{\min}\|w\|_{L^2(J;V)}^2.
\]
For such \(w\), \(\|w\|_X=\|w\|_{L^2(J;V)}\), while \(\|v\|_Y\le C\|w\|_X\) with \(C\) depending only on \(J\) and the embedding \(V\hookrightarrow H\).  These degenerate levels therefore also satisfy a uniform lower bound.  Consequently,
\[
  \inf_{L\in\Nzero}\gam(X_L,\hatY_L)>0 .
\]

\paragraph{Step 3: transfer from \(\hatY_L\) to \(Y_L\).}
It remains to prove the step omitted in the printed proof:
\[
  \inf_{L\in\Nzero}\K_{Y\times Y}(\hatY_L,Y_L)>0 .
\]
Set
\[
  W_0:=V_0,
  \qquad
  W_j:=V_j\cap V_{j-1}^{\perp_V}\quad (j\ge1).
\]
Then \(V_\ell=\bigoplus_{j=0}^\ell W_j\) orthogonally in \(V\).  By nestedness of the temporal spaces,
\[
  \hatY_{L,1}=\bigoplus_{j=0}^L \hatF_{L-j}\otimes W_j,
  \qquad
  Y_{L,1}=\bigoplus_{j=0}^L F_{L-j}\otimes W_j .
\]
Let \(\Pi_L:Y_1\to Y_{L,1}\) be the \(Y_1=L^2(J;V)\)-orthogonal projector.  For
\[
  z=\sum_{j=0}^L z_j,
  \qquad
  z_j\in\hatF_{L-j}\otimes W_j,
\]
the projector acts blockwise as
\[
  \Pi_L z=\sum_{j=0}^L (P_{L-j}\otimes I_{W_j})z_j .
\]
Using the lower bound from Step 1 and orthogonality of the \(W_j\)-blocks,
\[
\begin{aligned}
  \|\Pi_L z\|_{Y_1}^2
  &=\sum_{j=0}^L
    \|(P_{L-j}\otimes I_{W_j})z_j\|_{Y_1}^2  \\
  &\ge c_T^2\sum_{j=0}^L\|z_j\|_{Y_1}^2
   =c_T^2\|z\|_{Y_1}^2 .
\end{aligned}
\]
The projection characterization of \(K\) gives
\[
  \K_{Y_1\times Y_1}(\hatY_{L,1},Y_{L,1})\ge c_T .
\]
The second component is identical in \(\hatY_L\) and \(Y_L\), namely \(V_L\).  Hence
\[
  \K_{Y\times Y}(\hatY_L,Y_L)\ge c_T .
\]
Finally, the robustness estimate (4.4.34) gives
\[
  \gam(X_L,Y_L)
  \ge \K_{Y\times Y}(\hatY_L,Y_L)\,\gam(X_L,\hatY_L)
  \ge c_T\,\gam(X_L,\hatY_L).
\]
Taking the infimum over \(L\) proves the corrected theorem.

\section*{4. Variants in Remark 5.2.15}
For full tensor products, anisotropic sparse tensor products, and finite sums of sparse tensor products, the same proof applies after replacing the basic lower set
\[
  \{(k,\ell):0\le k+\ell\le L\}
\]
by the corresponding lower index set.  In Step 3, each spatial block \(W_j\) is paired with the largest active temporal level allowed by that lower set.  The same blockwise projection estimate then gives the identical lower bound \(c_T\).  This is the missing justification behind the ``without loss of generality'' reduction in the printed proof.

\end{document}
